\section{Related Work}\label{sec:related-work}
At the time of writing, the field of quantum error correction 
codes is experiencing a montage of \(70\) years worth of classical information theory, 
error correction codes and machine learning, making it hard to stay up to date.
In writing this section we aimed to give the reader a grounding in related QECC,
and specifically machine learning for ECC and QECC. 
Two resources were instrumental in laying the ground work in QECC for this work. 
The first is by Pantellev and Kalachev~\cite{panteleev2021degenerate}, 
which present BP+OSD decoding used here, as well as multiple code constructions, and methodology notes.
The second is Gottesman's book-in-writing~\cite{gottesman2026surviving}, 
which serves as a one-stop-shop for everything QECC related, 
and also makes a brilliant read\footnote{With pearls of wisdom such as ``...lattice surgery is an elective procedure...''}. 
The decoder implemented in this work is the one reported in~\cite{roffe_decoding_2020}, 
which is a CPU bound implementation. 
For the Belief Propagation step, specifically Min-Sum decoding, 
we considered using~\cite{sella2025mapping}, but have not yet paired it with the 
OSD post processing stage, which is needed to avoid CPU-GPU contention. 
For the BB family specifically, we note that search algorithms were reported by~\cite{wang2026coprime, wang2024coprime}.
\noindent Parametrised constructions of quantum Low Density Parity Check (qLDPC) Error Correcting Codes (ECC) 
draw from several mathematical constructs like polynomials~\cite{panteleev2021degenerate}
and graphs~\cite{evra2022decodable, manjunath2026universal}, 
but ultimately require an assignment of these parameters for practical use, 
and~\cite{panteleev2021degenerate} raises the question of finding good parameters. 

\noindent Reinforcement learning has been applied to the actual decoding as shown in~\cite{sweke2020reinforcement, blue2026machine}, as well as
to modify surface-code construction in~\cite{nautrup2019optimizing}, 
optimizing tensor-network codes as reported in~\cite{mauron2024optimization}, 
and reducing stabiliser weights of existing qLDPC codes~\cite{he2025discovering}. 
These works reward primarily structural properties of code parameters. 
In contrast, our main focus is the measured performance of the code under a fixed decoder.
Machine learning alternatives to RL have appeared recently as well, one workflow that uses LLM was introduced in~\cite{cruzbenito2026evolutionary}, and, closer to our \emph{encoder} approach, 
a Bayesian-optimisation framework that trains a model that predicts logical error rates without decoder simulation, 
and uses it to discover competitive \([[144,36]]\) and \([[144,16]]\) codes~\cite{cheng2026bayesian}. 
In our work this type of model is deployed to warm-start an RL agent, which trains a generative policy. 
Conceptually, the closest work to ours is~\cite{su2305discovery}, which sets to discover QECCs with \(k=1\) logical qubits using an analytic proxy physical-to-logical error rate. 
The work in~\cite{su2305discovery} applies RL methods to Quantum Lego~\cite{cao2022quantum}, a way to build QECCs from tensor networks.
Finally, this work is continuation of our previous work on co-designing (classical) LDPC codes and decoders~\cite{sella2024coding}. 